\section{Paradoksalny rozkład grupy wolnej}
\subsection{Podział według pierwszej litery}
Dla $s\in\mathcal A$ niech $W(s)$ oznacza zbiór niepustych słów zredukowanych zaczynających się literą $s$. Zatem
\begin{equation}\label{eq:4-1}
 F_2=\{e\}\sqcup W(a)\sqcup W(\bar a)\sqcup W(b)\sqcup W(\bar b).
\end{equation}
Każde słowo niepuste ma dokładnie jedną pierwszą literę, więc zbiory po prawej stronie są parami rozłączne i pokrywają $F_2$.
\begin{lemat}[Kasowanie pierwszej litery]\label{lem:pierwsza-litera}
Mamy $aW(\bar a)=F_2\setminus W(a)$ oraz $bW(\bar b)=F_2\setminus W(b)$.
\end{lemat}
\begin{proof}
Sprawdzimy obie inkluzje pierwszej równości. Słowo z $W(\bar a)$ ma postać $\bar au$, gdzie $u$ jest puste albo nie zaczyna się literą $a$; w przeciwnym razie nie byłoby zredukowane. Mnożenie przez $a$ z lewej kasuje parę $a\bar a$ i daje właśnie $u$. Wynik nie należy do $W(a)$.

Odwrotnie, jeśli $u\notin W(a)$, to $u=e$ albo jego pierwsza litera nie jest $a$. Słowo $\bar au$ jest wtedy zredukowane i należy do $W(\bar a)$, a $a(\bar au)=u$. Dotyczy to także $u=e$: jego poprzednikiem jest słowo $\bar a$. Drugą równość dowodzimy tak samo, zastępując $a$ przez $b$.
\end{proof}
Już teraz mamy dwa równania $F_2=W(a)\sqcup aW(\bar a)$ oraz $F_2=W(b)\sqcup bW(\bar b)$. Nie jest to jeszcze pełny rozkład paradoksalny, ponieważ cztery źródłowe zbiory $W(s)$ nie zawierają $e$. Nie wolno zgubić tego jednego elementu.

\subsection{Włączenie słowa pustego bez dodatkowej kopii}
Połóżmy $P=\{a^n:n=0,1,2,\ldots\}$, gdzie $a^0=e$. Zauważmy dokładnie, że
\begin{equation}\label{eq:4-2}
 aP=P\setminus\{e\}=P\cap W(a).
\end{equation}
Pierwsza równość wynika z $a\cdot a^n=a^{n+1}$ i z tego, że potęgi $a^0,a^1,\ldots$ są parami różne. Druga mówi, że spośród tych potęg tylko $a^0$ nie zaczyna się literą $a$.
\begin{twierdzenie}[Cztery kawałki grupy]\label{tw:cztery-kawalki}
Zbiory
\begin{equation}\label{eq:4-3}
 A_1=W(a)\setminus P,\quad A_2=W(\bar a)\cup P,
 \qquad B_1=W(b),\quad B_2=W(\bar b)
\end{equation}
tworzą partycję $F_2$ i spełniają
\begin{equation}\label{eq:4-4}
 F_2=A_1\sqcup aA_2,\qquad F_2=B_1\sqcup bB_2.
\end{equation}
W szczególności $F_2$ jest zbiorem paradoksalnym względem działania na sobie przez mnożenie z lewej.
\end{twierdzenie}
\begin{proof}
\emph{Partycja źródła.} W partycji \eqref{eq:4-1} wydzielamy dodatnie potęgi $a$ ze zbioru $W(a)$ i razem z $e$ dołączamy je do $W(\bar a)$, tworząc $A_2=W(\bar a)\cup P$. Pozostała część $W(a)$ to $A_1$. Nic nie zmieniamy w $W(b)$ i $W(\bar b)$. Zbiory pozostają rozłączne, a teraz pokrywają także $e$.

\emph{Pierwsza kopia.} Korzystając z lematu~\ref{lem:pierwsza-litera} i \eqref{eq:4-2}, otrzymujemy $aA_2=aW(\bar a)\sqcup aP=(F_2\setminus W(a))\sqcup(P\setminus\{e\})$. Pierwszy składnik nie zawiera słów zaczynających się od $a$, a drugi składa się właśnie z dodatnich potęg $a$. Ich unia jest rozłączna. Zbiór $A_1=W(a)\setminus P$ jest rozłączny z obydwoma składnikami i uzupełnia je do całego $F_2$.

\emph{Druga kopia.} Lemat~\ref{lem:pierwsza-litera} od razu daje $B_1\sqcup bB_2=W(b)\sqcup(F_2\setminus W(b))=F_2$. W obu kopiach sprawdziliśmy zarówno pokrycie, jak i rozłączność.
\end{proof}
\begin{center}
\begin{tabular}{c|c|c}
 kawałek źródłowy&ruch&etykieta celu\\\hline
 $A_1$&$e$&$1$\\ $A_2$&$a$&$1$\\ $B_1$&$e$&$2$\\ $B_2$&$b$&$2$
\end{tabular}
\end{center}
\begin{przyklad}[Sprawdzenie kilku punktów]
$e,a,a^2$ należą do $A_2$ i po ruchu $a$ trafiają odpowiednio do $a,a^2,a^3$ w pierwszej kopii. Element $\bar a\in A_2$ trafia do $e$ w pierwszej kopii. Słowo $ab\in A_1$ pozostaje $ab$ w tej samej kopii. Element $\bar b\in B_2$ trafia do $e$ w drugiej kopii. Dwie kopie $e$ w zbiorze docelowym mają różnych poprzedników i różne etykiety.
\end{przyklad}

